
\RequirePackage{fix-cm}
\documentclass[final]{article}

% =========================================================
% FORMATO GENERAL: A1 VERTICAL
% =========================================================
\usepackage[
    paperwidth=59.4cm,
    paperheight=84.1cm,
    margin=0.75cm
]{geometry}
\pagestyle{empty}

\usepackage[utf8]{inputenc}
\usepackage[T1]{fontenc}
\usepackage[spanish]{babel}
\usepackage{anyfontsize}

\usepackage{graphicx}
\usepackage{xcolor}
\usepackage{booktabs}
\usepackage{tabularx}
\usepackage{array}
\usepackage{tikz}
\usepackage[most]{tcolorbox}
\usepackage{caption}
\usepackage{enumitem}
\usepackage{amsmath}
\usepackage{ragged2e}
\usepackage{multicol}

% =========================================================
% COLORES
% =========================================================
\definecolor{mainblue}{HTML}{12355B}
\definecolor{accent}{HTML}{C0392B}
\definecolor{softblue}{HTML}{EAF2F8}
\definecolor{softgray}{HTML}{F7F7F7}
\definecolor{darkgray}{HTML}{2E2E2E}
\definecolor{goodgreen}{HTML}{1B7F5C}
\definecolor{orange}{HTML}{D68910}

% =========================================================
% ESTILO GENERAL
% =========================================================
\setlength{\parindent}{0pt}
\setlength{\parskip}{0.17em}
\setlist[itemize]{
    leftmargin=1.12em,
    itemsep=0.04em,
    topsep=0.04em,
    parsep=0pt,
    partopsep=0pt
}

\captionsetup{
    font=small,
    labelfont=bf,
    justification=centering,
    skip=2pt
}

\renewcommand{\arraystretch}{1.08}

% =========================================================
% COMANDOS
% =========================================================
\newcommand{\PosterTitle}[1]{%
    {\fontsize{50}{56}\selectfont\bfseries\color{white} #1}
}

\newcommand{\PosterSubtitle}[1]{%
    {\fontsize{19.5}{23}\selectfont\color{white} #1}
}

\newcommand{\sectionbox}[2]{%
\begin{tcolorbox}[
    enhanced,
    colback=white,
    colframe=mainblue,
    boxrule=1.22pt,
    arc=2.5mm,
    left=3mm,
    right=3mm,
    top=2.2mm,
    bottom=2.2mm,
    title=\fontsize{19.4}{22.8}\selectfont\bfseries #1,
    coltitle=white,
    colbacktitle=mainblue,
    fonttitle=\bfseries,
    attach boxed title to top left={xshift=3mm,yshift=-1.55mm},
    boxed title style={arc=1.8mm,boxrule=0pt},
    fontupper=\fontsize{15.1}{17.6}\selectfont
]
#2
\end{tcolorbox}
\vspace{0.04cm}
}

\newcommand{\metricbox}[3]{%
\begin{tcolorbox}[
    enhanced,
    colback=#1!8,
    colframe=#1,
    boxrule=0.9pt,
    arc=1.8mm,
    left=1mm,
    right=1mm,
    top=0.9mm,
    bottom=0.9mm
]
\centering
{\fontsize{25}{28}\selectfont\bfseries #2}\\[-0.15em]
{\fontsize{10.8}{12.3}\selectfont #3}
\end{tcolorbox}
}

\newcommand{\posterimage}[5][]{%
{\centering
    \includegraphics[
        width=\linewidth,
        height=#3,
        keepaspectratio,
        #1
    ]{#2}\par
    \vspace{-0.15cm}
    \captionof{figure}{#4}
    \vspace{#5}
}
}

\begin{document}

% Todo el póster en un único bloque para evitar que el título quede separado.
\noindent
\begin{minipage}[t]{\textwidth}

% =========================================================
% ENCABEZADO
% =========================================================
\begin{tcolorbox}[
    enhanced,
    colback=mainblue,
    colframe=mainblue,
    arc=0mm,
    boxrule=0pt,
    left=8mm,
    right=8mm,
    top=4.4mm,
    bottom=4.4mm
]
\PosterTitle{Clasificación y reconstrucción de objetos en visión biónica}\\[0.18em]
\PosterSubtitle{Evaluación de información semántica y espacial preservada en perceptos simulados con \texttt{pulse2percept}}\\[0.20em]
{\fontsize{18}{21}\selectfont\color{white}
\textbf{Joel Jablonski, Estanislao Molina Abeniacar}
\hfill
Universidad de San Andrés -- Visión Artificial
}
\end{tcolorbox}

\vspace{0.20cm}

% =========================================================
% DOS COLUMNAS NORMALES
% =========================================================
\begin{multicols}{2}
\setlength{\columnsep}{0.9cm}

% ---------------------------------------------------------
\sectionbox{1. Problema y motivación}{
Las prótesis retinales no transmiten una imagen natural completa: generan una percepción degradada, de baja resolución y dependiente de la geometría del implante.

\textbf{Motivación real.}
El objetivo de fondo es estudiar qué transformaciones o filtros podrían hacer que una escena sea más interpretable para una persona usuaria de una prótesis visual.

Como evaluar filtros directamente con usuarios requiere experimentos psicofísicos y métricas perceptuales, este trabajo aborda una etapa previa:

\begin{center}
\textbf{medir cuánta información útil sobrevive al pasaje por un implante simulado.}
\end{center}

Para eso se usaron dos tareas proxy:
\begin{itemize}
    \item \textbf{Clasificación}: mide preservación semántica.
    \item \textbf{Reconstrucción}: mide preservación espacial y de forma.
\end{itemize}
}

% ---------------------------------------------------------
\sectionbox{2. Pipeline general}{
\posterimage[]
{figs/FOTO1.pdf}
{12cm}
{Pipeline general: objetos segmentados de COCO se transforman en perceptos simulados y luego se evalúan mediante clasificación y reconstrucción.}
{0cm}

El flujo experimental fue:
\[
\text{COCO} \rightarrow \text{máscara} \rightarrow \texttt{pulse2percept} \rightarrow \text{percepto}
\]

Luego, cada percepto se utilizó como entrada para:
\[
\text{percepto} \rightarrow \text{clase}
\qquad
\text{y}
\qquad
\text{percepto} \rightarrow \text{máscara reconstruida}
\]
}

% ---------------------------------------------------------
\sectionbox{3. Conjunto de datos y perceptos}{
Se trabajó con máscaras binarias de COCO para cuatro clases:

\begin{center}
\textbf{person} \quad
\textbf{car} \quad
\textbf{chair} \quad
\textbf{dining table}
\end{center}

Las máscaras se normalizaron a $256 \times 256$ píxeles manteniendo aspect ratio mediante padding negro.

\begin{center}
{\fontsize{11.8}{13.5}\selectfont
\begin{tabular}{lc}
\toprule
\textbf{Clase} & \textbf{Muestras} \\
\midrule
person & 900 \\
car & 900 \\
chair & 900 \\
dining table & 624 \\
\bottomrule
\end{tabular}
}
\end{center}

Se generaron perceptos para tres modelos de implante:
\begin{itemize}
    \item \textbf{Argus II}
    \item \textbf{PRIMA}
    \item \textbf{Alpha-AMS}
\end{itemize}

La partición train/validación/test se realizó por \texttt{sample\_id}, no por fila, para evitar que perceptos del mismo objeto aparezcan en splits distintos.
}

% ---------------------------------------------------------
\sectionbox{4. Ejemplos de perceptos}{
\posterimage[]
{figs/FOTO2.PNG}
{13cm}
{Ejemplos de perceptos simulados. Una misma máscara genera representaciones distintas según el implante.}
{0cm}

\texttt{pulse2percept} se utilizó como herramienta de simulación controlada: no modela toda la experiencia visual de un paciente, sino la degradación perceptual asociada a distintas geometrías de implante.

Esto permite comparar si distintas representaciones conservan información suficiente para reconocer o reconstruir el objeto original.
}

% ---------------------------------------------------------
\sectionbox{5. Clasificación de objetos}{
\textbf{Entrada:} percepto simulado. \\
\textbf{Salida:} clase del objeto.

Se compararon modelos clásicos y neuronales:
\begin{itemize}
    \item \textbf{KNN + PCA}: percepto redimensionado, aplanado y clasificado por similitud.
    \item \textbf{CNN}: modelo neuronal por implante individual.
    \item \textbf{CNN + GRU}: fusión aprendida de perceptos multi-implante.
\end{itemize}

\posterimage[]
{figs/fig4.pdf}
{20cm}
{Arquitecturas de clasificación: CNN individual y CNN+GRU para fusión multi-implante.}
{0cm}
}

\columnbreak

% -------------------------------------------------------
\sectionbox{6. Resultados de clasificación}{
\centering
{\fontsize{12.5}{14.2}\selectfont
\begin{tabularx}{\linewidth}{>{\raggedright\arraybackslash}X c}
\toprule
\textbf{Modelo} & \textbf{Accuracy test} \\
\midrule
Baseline mayoría & 0.271 \\
KNN multi-implante & \textbf{0.611} \\
CNN + GRU & 0.575 \\
KNN PRIMA & 0.571 \\
KNN Alpha-AMS & 0.567 \\
CNN PRIMA & 0.563 \\
KNN Argus II & 0.561 \\
CNN Alpha-AMS & 0.561 \\
CNN Argus II & 0.509 \\
\bottomrule
\end{tabularx}
}

\vspace{0.18em}

\begin{minipage}{0.32\linewidth}
\metricbox{goodgreen}{0.611}{KNN\\multi-implante}
\end{minipage}
\hfill
\begin{minipage}{0.32\linewidth}
\metricbox{orange}{0.575}{CNN+GRU\\mejor red}
\end{minipage}
\hfill
\begin{minipage}{0.32\linewidth}
\metricbox{mainblue}{0.271}{baseline\\mayoría}
\end{minipage}

\vspace{0.10cm}

\posterimage[]
{figs/confusion_knn_multi.PNG}
{15cm}
{Matriz de confusión del mejor clasificador: KNN multi-implante.}
{0cm}

\justifying
El resultado es contraintuitivo: el modelo clásico KNN multi-implante superó a CNN+GRU. Esto sugiere que, para este tamaño de dataset, PCA + similitud directa fue más estable que entrenar una red neuronal desde cero.
}

% ---------------------------------------------------------
\sectionbox{7. Reconstrucción de máscaras}{
\textbf{Entrada:} percepto individual o perceptos fusionados. \\
\textbf{Salida:} máscara binaria reconstruida.

Se utilizó una arquitectura \textbf{U-Net}, adecuada para predicción píxel a píxel:
\begin{itemize}
    \item el encoder extrae estructura global;
    \item el decoder recupera resolución espacial;
    \item las skip connections preservan localización y bordes.
\end{itemize}

\posterimage[]
{figs/arquitectura_unet.pdf}
{15cm}
{U-Net de fusión: los perceptos de los tres implantes se apilan como canales de entrada y la salida es una máscara binaria.}
{0cm}
}

% ---------------------------------------------------------
\sectionbox{8. Resultados de reconstrucción}{
\centering
{\fontsize{12.3}{14}\selectfont
\begin{tabularx}{\linewidth}{>{\raggedright\arraybackslash}X c c c}
\toprule
\textbf{Modelo} & \textbf{Dice} & \textbf{IoU} & \textbf{Acc. píxel} \\
\midrule
U-Net Argus II & 0.883 & 0.835 & 0.986 \\
U-Net PRIMA & 0.871 & 0.825 & 0.982 \\
U-Net Alpha-AMS & 0.902 & 0.862 & 0.986 \\
U-Net fusión & \textbf{0.905} & \textbf{0.876} & \textbf{0.989} \\
\bottomrule
\end{tabularx}
}

\vspace{0.18em}

\begin{minipage}{0.32\linewidth}
\metricbox{goodgreen}{0.905}{Dice\\U-Net fusión}
\end{minipage}
\hfill
\begin{minipage}{0.32\linewidth}
\metricbox{goodgreen}{0.876}{IoU\\U-Net fusión}
\end{minipage}
\hfill
\begin{minipage}{0.32\linewidth}
\metricbox{goodgreen}{0.989}{Acc. píxel\\U-Net fusión}
\end{minipage}

\vspace{0.10cm}

\posterimage[]
{figs/FOTO6.PNG}
{9.1cm}
{Reconstrucciones cualitativas: la U-Net recupera la forma global, aunque pierde bordes finos y partes pequeñas.}
{0cm}

\justifying
La reconstrucción fue más sólida que la clasificación. Sin embargo, la pixel accuracy debe interpretarse con cuidado, porque gran parte de la imagen corresponde al fondo negro. Dice e IoU son más informativas para evaluar superposición de la región del objeto.
}

% ---------------------------------------------------------
\sectionbox{9. Conclusiones e interpretación}{
\begin{itemize}
    \item Los perceptos simulados conservan información útil, aunque degradada.
    \item La \textbf{fusión multi-implante} fue el mejor enfoque tanto para clasificación como para reconstrucción.
    \item La U-Net reconstruyó correctamente la \textbf{silueta global} de los objetos, pero perdió detalles finos.
    \item El KNN multi-implante superó a CNN+GRU, sugiriendo que los patrones globales de forma e intensidad fueron discriminativos.
    \item El trabajo no constituye una validación clínica, sino una etapa previa para comparar representaciones visuales en visión protésica simulada.
\end{itemize}
}

% ---------------------------------------------------------
\sectionbox{10. Uso real y trabajo futuro}{
Este pipeline puede servir para evaluar filtros candidatos antes de realizar pruebas con usuarios:
\[
\text{filtro candidato}
\rightarrow
\text{percepto simulado}
\rightarrow
\text{clasificación / reconstrucción}
\]

Si un filtro mejora la accuracy o el Dice, entonces preserva más información semántica o espacial.

\textbf{Trabajo futuro:}
\begin{itemize}
    \item comparar filtros visuales simples: bordes, segmentación, siluetas, saliency;
    \item extender a imágenes RGB y escenas completas;
    \item evaluar secuencias de video;
    \item entrenar modelos surrogate diferenciables para optimización end-to-end.
\end{itemize}
}

\end{multicols}

% =========================================================
% PIE
% =========================================================
\begin{tcolorbox}[
    enhanced,
    colback=softgray,
    colframe=mainblue,
    boxrule=0.75pt,
    arc=1.5mm,
    left=3mm,
    right=3mm,
    top=1.2mm,
    bottom=1.2mm
]
{\fontsize{9.2}{10.8}\selectfont
\textbf{Referencias clave:}
Beyeler et al., \emph{pulse2percept: A Python-Based Simulation Framework for Bionic Vision};
Han et al., \emph{Deep Learning-Based Scene Simplification for Bionic Vision};
Sanchez-Garcia et al., \emph{Semantic and Structural Image Segmentation for Prosthetic Vision};
Lin et al., \emph{Microsoft COCO: Common Objects in Context}.
\hfill
\\
\textbf{Pipeline:} COCO $\rightarrow$ máscaras $\rightarrow$ \texttt{pulse2percept} $\rightarrow$ clasificación y reconstrucción.
}
\end{tcolorbox}

\end{minipage}

\end{document}
